\documentclass[11pt]{article}
\usepackage[margin=1in]{geometry}\usepackage{amsmath,amssymb}
\title{A one-bit condition can leave the semisimple-dual family}
\date{28 October 2026}\begin{document}\maketitle
This obstruction concerns closure of the inherited certificate, not
hardness of the resulting count or rejection by every possible frontend.
Work over $\mathbb F_2$ in $M_8(\mathbb F_2)$ and put
\[
 \mathcal D=\{\operatorname{diag}(A,A,B,B):A,B\in M_2(\mathbb F_2)\},
 \qquad W=\mathcal D^\perp
\]
under $\langle D,Y\rangle=\operatorname{tr}(DY)$.
The right entry $Y_{31}$ is an independent free coordinate of $W$.
Pinning that one input bit to zero changes the reverse dual to
\[
 \mathcal D'=\mathcal D+\langle E_{13}\rangle,
\]
since $\operatorname{tr}(E_{13}Y)=Y_{31}$ and $E_{13}\notin\mathcal D$.

Its actual left multiplier algebra is
\[
 \{T:T\mathcal D'\subseteq\mathcal D'\}
 =\{\operatorname{diag}(A,A,B,B): A_{21}=0\}
 \simeq \operatorname{UT}_2(\mathbb F_2)\oplus M_2(\mathbb F_2).
\]
Indeed $I\in\mathcal D'$, so any multiplier has form
$T=\operatorname{diag}(A,A,B,B)+bE_{13}$. Multiplication by an old
diagonal element with first block $C$ creates the term $bC_{12}E_{14}$,
forcing $b=0$. Then $TE_{13}=A_{11}E_{13}+A_{21}E_{23}$, forcing
$A_{21}=0$. Conversely these conditions suffice on every basis element.

The first factor has the nonzero square-zero radical generated by
$\operatorname{diag}(E_{12},E_{12},0,0)$. Therefore the retained full
simple-factor/sandwich test cannot certify this multiplier algebra as
semisimple. The regular two-factor representation is not closed under
even this legal one-bit restriction. A general prefix sampler cannot
be inferred from the count theorem without an additional representation
or direct fiber-sampling argument. Signed histogram coefficients likewise
do not provide nonnegative branching probabilities.
\end{document}
